% from ega0-22-fr.tex
\oldpage[0\textsubscript{IV-1}]{186}

Or, la formation par récurrence de $v_j$ résulte de l'hypothèse sur $B$ et
du même raisonnement que dans (19.3.10.1), en utilisant le lemme
(19.3.10.2) et le fait que
\[
  \mathfrak n'^j+(\mathfrak n'^{j+1}+\mathfrak mB')
  =\mathfrak n'^j+\mathfrak mB'.
\]
Il résulte alors de (19.7.1.4) que $v$ est un $A$-\emph{isomorphisme}, car
$B'$ est un $A$-module plat, et $B$ est complet pour la topologie
$\mathfrak m$-préadique, les idéaux $\mathfrak m^iB$ étant fermés dans $B$
pour la topologie $\mathfrak n$-adique ($\mathbf 0_{\mathrm I}$, 7.3.5).
C.Q.F.D.
\end{proof}

\subsection{Caractérisation différentielle des algèbres locales formellement
lisses sur un corps}
\label{subsection:0.22.2-fr}

\begin{env}[22.2.1]
\label{0.22.2.1-fr}
Soient $k$ un corps, $P$ son sous-corps premier, $A$ une $k$-algèbre qui est
un anneau \emph{local}, $\mathfrak m$ son idéal maximal, $K=A/\mathfrak m$
son corps résiduel. Comme $K$ est séparable sur $P$, $K$ est une
$P$-algèbre formellement lisse pour la topologie discrète (19.6.1), et par
suite la $P$-extension $A/\mathfrak m^2$ de $K$ par
$\mathfrak m/\mathfrak m^2$ est $P$-\emph{triviale}, et l'homomorphisme
caractéristique
\begin{equation}
\tag{22.2.1.1}
\label{0.22.2.1.1-fr}
  \chi_{A/k}:\Upsilon_{K/k}\longrightarrow\mathfrak m/\mathfrak m^2
\end{equation}
est donc défini (20.6.23).
\end{env}

\begin{theorem}[22.2.2]
\label{0.22.2.2-fr}
Sous les conditions de (22.2.1), pour que $A$ soit une $k$-algèbre
formellement lisse pour la topologie $\mathfrak m$-préadique, il faut et il
suffit que les conditions suivantes soient satisfaites~:
\begin{enumerate}
  \item[(i)] L'homomorphisme canonique
  $\mathbf S_K^*(\mathfrak m/\mathfrak m^2)\to
  \operatorname{gr}_{\mathfrak m}^*(A)$ est bijectif.
  \item[(ii)] L'homomorphisme caractéristique
  $\chi_{A/k}:\Upsilon_{K/k}\to\mathfrak m/\mathfrak m^2$ est injectif.
\end{enumerate}
\end{theorem}

\begin{proof}
C'est un cas particulier de (22.1.2), où l'on remplace $\Lambda$ par $P$,
$A$ par $k$ et $B$ par $A$~; $K$ et $k$, étant séparables sur $P$, sont en
effet des $P$-algèbres formellement lisses (19.6.1).
\end{proof}

\begin{remark}[22.2.3]
\label{0.22.2.3-fr}
La condition (ii) de (22.2.2) peut être remplacée par une quelconque des
suivantes~:
\begin{enumerate}
  \item[a)] L'homomorphisme canonique
  $\Omega^1_k\otimes_kK\to\Omega^1_A\otimes_AK$ est injectif.
  \item[b)] Pour tout $n$, l'homomorphisme canonique
  \[
    \Omega^1_k\otimes_k(A/\mathfrak m^{n+1})
    \longrightarrow
    \Omega^1_A\otimes_A(A/\mathfrak m^{n+1})
  \]
  est inversible à gauche.
  \item[c)] L'homomorphisme canonique
  $\Omega^1_k\widehat\otimes_kA\to\widehat\Omega^1_A$ est inversible à
  gauche.
  \item[d)] $A$ est une $k$-algèbre formellement lisse (pour la topologie
  $\mathfrak m$-préadique) \emph{relativement au corps premier}.
\end{enumerate}
En effet, on a déjà remarqué dans (22.1.1) que l'injectivité de $\chi_{A/k}$
est équivalente à $a)$, et que la conjonction de $a)$ et de la condition (i)
de (22.2.2) équivaut à celle de $d)$ et de (i). Enfin, on a aussi vu dans
(22.1.1) que $b)$ et $d)$ sont équivalentes, et l'équivalence de $b)$ et
$c)$ résulte de (19.1.8).
\end{remark}
\begin{env}[22.2.4]
\label{0.22.2.4-fr}
La principale application de (22.2.2) a trait aux $k$-algèbres locales
noethériennes et \emph{complètes} (cas où $\mathfrak m/\mathfrak m^2$ est un
$K$-espace vectoriel de rang \emph{fini}). De façon précise, soient $k$ un
corps, $K$ une extension de $k$, $V$ un $K$-espace vectoriel de

